\label{ch:gravity}

\section{Constitutive chain of the gravitational sector: chronological scale \(108\), tensor reduction, holonomy and metrological evaluation of \(G\)}

Gravity is organized here into four strata that must not be merged:

\begin{equation}
\label{eq:gravity-causal-chain}
\begin{aligned}
\text{CT108 chronological structure and action lines}
&\longrightarrow \text{periodicity selector }\theta_{108}^{\rm circ}
\longrightarrow G_{108},\\
\text{icosahedral incidence}
&\longrightarrow V_{12}\xrightarrow{P_5}V_5
\xrightarrow{\Gamma_5}\operatorname{STF}_3
\xrightarrow{\Lambda(n)}\mathcal H_{\rm TT}(n),\\
\text{enriched holonomic route}
&\dashrightarrow G_{\rm hol},\\
\text{metrological chart}
&\longrightarrow 6.67430\times10^{-11}\,
\mathrm{m^3\,kg^{-1}\,s^{-2}}.
\end{aligned}
\end{equation}

The first line states a condition of periodicity that is intrinsic and exact.
The second constructs the carrier and its trace-free transverse reduction. The third preserves an
explicit construction boundary. The fourth is an evaluation in a
declared decimal chart: it identifies the magnitude but does not select the gravitational
character.

\section{Gravitational magnitude line and chronological scale \(108\)}

Let \(c_{\rm int}>0\), \(t_0>0\),
\(\ell_0=c_{\rm int}t_0\) and a positive section
\(\hbar_a\) of the action line, with
\(a\in\{\mathrm{pre},\mathrm{ret}\}\). The dimensional transducer is

\begin{equation}
\label{eq:gravity-transducer}
G_a(\theta)
=\theta\frac{c_{\rm int}^{5}t_0^2}{\hbar_a}
=\theta\frac{c_{\rm int}^{3}\ell_0^2}{\hbar_a},
\qquad \theta>0.
\end{equation}

The equality of the two expressions uses only
\(\ell_0=c_{\rm int}t_0\), and
\([G]=L^3M^{-1}T^{-2}\). The dimension is fixed by the line; the
mathematical problem is to select the dimensionless character \(\theta\).

The CT108 calendar defines

\begin{equation}
\label{eq:gravity-clock108}
T_{108}=108t_0,
\qquad
\omega_{108}=\frac{2\pi}{108t_0},
\qquad
R_{108}=\frac{c_{\rm int}}{\omega_{108}}
=\frac{54}{\pi}\ell_0.
\end{equation}

For each action sheet, let

\begin{equation}
\label{eq:gravity-clock-energy-mass}
E_{108,a}=\hbar_a\omega_{108},
\qquad
m_{108,a}=\frac{E_{108,a}}{c_{\rm int}^2}.
\end{equation}

Then, the reduced Compton wavelength coincides with the radius of the cycle:

\begin{equation}
\label{eq:gravity-compton-clock}
\bar\lambda_{C,a}
=\frac{\hbar_a}{m_{108,a}c_{\rm int}}
=R_{108}.
\end{equation}

\begin{theorem}[Uniqueness of the selector under the HMT periodicity condition]
\label{thm:gravity-circular-selector}
With the reduced convention
\(r_g=Gm/c_{\rm int}^2\) and the HMT periodicity premise that identifies the gravitational
radius of the CT108 quantum with its phase radius, the closure condition
\begin{equation}
\label{eq:gravity-closing-condition}
r_{g,a}(\theta)=R_{108}
\end{equation}
selects a unique positive solution,
\begin{equation}
\label{eq:gravity-theta108}
\boxed{
\theta_{108}^{\rm circ}=\left(\frac{54}{\pi}\right)^2
=295.4525715010569415303525\ldots .}
\end{equation}
In that embodiment,
\begin{equation}
\label{eq:gravity-four-lengths}
R_{108}=\bar\lambda_{C,a}=r_{g,a}=\ell_P^{\rm circ}.
\end{equation}
\end{theorem}

\begin{proof}
By \cref{eq:gravity-transducer,eq:gravity-clock-energy-mass},
\[
r_{g,a}(\theta)
=\frac{G_a(\theta)m_{108,a}}{c_{\rm int}^2}
=\theta c_{\rm int}t_0^2\omega_{108}
=\theta\frac{\pi}{54}\ell_0.
\]
Comparing it with
\(R_{108}=(54/\pi)\ell_0\) necessarily gives
\(\theta=(54/\pi)^2\). The equation is linear in \(\theta\) and all the
factors are positive, so the solution is unique.
\end{proof}

The equality \(r_g=R_{108}\) does not follow from the dimensionality of \cref{eq:gravity-transducer}. It is the explicit geometric condition of this HMT realization: action, Compton, and periodic return must determine the same length. The theorem proves that, once this structural condition is accepted, the character \(\theta\) is fixed without using the SI value of \(G\).

\begin{remark}[Radius convention]
The theorem uses the reduced gravitational radius \(Gm/c^2\). If it is replaced
by the Schwarzschild radius \(2Gm/c^2\), the selector changes by a factor of
two. The two expressions correspond to distinct geometric conditions,
so the convention must be declared before the calculation.
\end{remark}

\section{Reciprocity \(G\leftrightarrow\hbar\) and conservation of the Planck area}

\begin{corollary}[Planck family associated with the temporal periodicity of \(108\) phases]
\label{cor:gravity-planck-family}
For \(G_a=G_a(\theta_{108}^{\rm circ})\),
\begin{align}
\ell_P&=\frac{54}{\pi}\ell_0,
&t_P&=\frac{54}{\pi}t_0,
\label{eq:gravity-planck-length-time}\\
E_{P,a}&=\frac{\hbar_a}{t_P}
=\frac{\pi\hbar_a}{54t_0}
=\frac{h_a}{108t_0},
&\ell_P^2&=\frac{G_a\hbar_a}{c_{\rm int}^3}
=\theta_{108}^{\rm circ}\ell_0^2.
\label{eq:gravity-planck-energy-area}
\end{align}
Adem\'as,
\begin{equation}
\label{eq:gravity-planck-force-power-density}
F_{P,a}=\frac{c_{\rm int}^4}{G_a},
\qquad
P_{P,a}=\frac{c_{\rm int}^5}{G_a},
\qquad
\rho_{P,a}=\frac{\hbar_a}
{(\theta_{108}^{\rm circ})^2c_{\rm int}^5t_0^4}.
\end{equation}
\end{corollary}

\begin{proof}
The first three identities follow from
\(\ell_P=R_{108}\), \(t_P=\ell_P/c_{\rm int}\), and
\(h_a=2\pi\hbar_a\). For the area,
\[
\frac{G_a\hbar_a}{c_{\rm int}^3}
=\theta_{108}^{\rm circ}c_{\rm int}^2t_0^2
=\theta_{108}^{\rm circ}\ell_0^2.
\]
The remaining expressions are the dimensional compositions of force,
power, and mass per volume with the same section.
\end{proof}

Let
\begin{equation}
\label{eq:gravity-Ract}
R_{\rm act}:=\frac{\hbar_{\rm pre}}{\hbar_{\rm ret}}.
\end{equation}
Since the action appears in the numerator of the chronological mass and in the
denominator of \(G\),
\begin{equation}
\label{eq:gravity-twoway-covariance}
\frac{m_{108,\rm pre}}{m_{108,\rm ret}}=R_{\rm act},
\qquad
\frac{G_{\rm pre}}{G_{\rm ret}}=R_{\rm act}^{-1},
\qquad
G_a\hbar_a
=\theta_{108}^{\rm circ}c_{\rm int}^5t_0^2.
\end{equation}
Therefore, \(\ell_P^2=G_a\hbar_a/c^3\) does not change sheets. Gravity and
action are reciprocally oriented, while area preserves the
common geometry. This cancellation is the interface of confluence toward the area
operator of the \cref{ch:modular}.
